\scanpage{29}
\[
\deleted{\unsure{\varphi}}\Sigma^n(\alpha_{s'})=\Sigma^n(\alpha_s)\underbrace{\Sigma^n(\operatorname{trans}\tau)}_{\varphi_n(\tau)} \tag*{(31)}
\]
et \deleted{\unsure{on a}} on a
\[
\begin{aligned}
&\varphi_n(\tau):\ \Sigma^n\check{\mathbb V}L)\xrightarrow{\sim}\Sigma^n\check{\mathbb V}L\\
&\sum^{\prime\prime} L^{\otimes i}\mathrel{\unsure{\longleftarrow|\longrightarrow}}\sum^{\prime\prime}L^{\otimes i}
\end{aligned}
\]
défini par la condition que l'on a commutativité de

\diagram{\begin{tikzcd}[ampersand replacement=\&,column sep=large,row sep=small]
\deleted{\Sigma^n\check{\mathbb V}L} \& {}\\
\check{\mathbb V}\sum_{1\le i\le n}L^{\otimes i} \arrow[r,"\varphi_n(\tau)"] \& \check{\mathbb V}\sum_{1\le i\le n}L^{\otimes i}\\
\check{\mathbb V}L^n \arrow[u,"\sigma_n"] \arrow[r,"(\operatorname{trans}\tau)^n"] \& \check{\mathbb V}(L^n) \arrow[u,"\sigma_n"]
\end{tikzcd}}

i.e.
\[
\left(\sum_{1\le i\le n}\deleted{\unsure{\zeta}}(x_1+\tau)(x_{\unsure{2|1}}\unsure{-|+}\tau)\text{---}(x_i+\tau)\right)
=\varphi_n(\tau)\underbrace{\left(\sum_i x_1\text{---}x_i\right)_i}_{\sigma_*(x)}
\]
\[
\begin{aligned}
&\Vert\\[-2pt]
&\deleted{\sigma_{\unsure{1|0}}(x)+n\tau,\ \sigma_2(x)+(n-1)\sigma_1(x)+\tfrac{n(n-1)}2\tau^2,\ \sigma_3(x)+}\\
&\deleted{\unsure{\sum}}\left(\prod_{1\le i\le n}(1+\tau\deleted{\unsure{i}}+x_i)\right)-1
=\left(\sum_{\unsure{0|1}\le p\le n}(1+\tau)^{n-p}\sigma_p\right)-1\\
&=\deleted{\sum_{\unsure{0}\le p}(1+\tau)^{\unsure{n}}}
\sum_{\substack{0\le p,q\le n\\p+q>0}}\frac{(n-p)!}{q!(n-p-q)!}\tau^q\sigma_p
\end{aligned}
\]
\margin{\uline{NB} On pose $\sigma_0=1$}

donc
\[
\varphi_n(\tau)(\sigma_1,\text{---},\sigma_n)=
 \bigl(\varphi_n(\tau)_i(\sigma_1,\text{---},\sigma_n)\bigr)_{1\le i\le n} \tag*{(32)}
\]
avec
\[
\varphi_n(\tau)_i(\sigma_1,\text{---},\sigma_n)=
\underbrace{\sum_{\substack{\deleted{\unsure{0\le}}p,q\ge0\\p+q=i\\\deleted{\unsure{p+q>0}}}}
\frac{(n-p)!}{(n-i\deleted{\unsure{-p}})!\,q!}\tau^q\sigma_p}_{\displaystyle =\ \deleted{\unsure{\varphi_i(\tau)(\sigma_1,\text{---},\sigma_n)}}}
\tag*{(33)}
\]
On voit que les $\varphi_n(\tau)$ sont des transformations affines

% archive sheet number: 28

